\begin{frame}[allowframebreaks]{Action potential(Spike)}
    When a stimulus reaches the threshold, a neuron depolarizes, causing a rapid change in the electrical potential difference across its membrane \cite{grider-2023-physiology-action-potential}.
    
    
    
    Similarly when a sensor is provoked by a stimulus, then a voltage difference between the nucleus and axon terminals is created which is called \textbf{action potential} \cite{hodgkin-1952-a-quantitative-description-of-membrane-current-and-its-application-to-conduction-and-excitation-in-nerve}
    \begin{itemize}
        \item Action potential has an oscillatory behavior. That is, the volatage goes first high and then goes lower than zero and then returns back to zero
    \end{itemize}
    \begin{figure}[H]
        \centering
        \includegraphics[width=0.5\textwidth]{/home/donkarlo/Dropbox/repo/nd_research_publication/src/kind/journal_2026/presentation/anatomy/body/nervous_system/neuron/action_potential/wave_form.png}
    \end{figure}
\end{frame}